\section{Behaviour as \texorpdfstring{$p\to\tfrac12^-$}{p to 1/2}}
\label{a:sec:nearcritical}

The bounds \cref{a:eq:Abounds} show that $A(p)$ is of order $\varepsilon$ near
criticality but leave a factor of two unresolved. The series \cref{a:eq:Aseries}
settles it. It also identifies the order of the first relative correction,
which turns out to carry a logarithm and is the reason the leading asymptotic is a
poor surrogate except very close to $p=\tfrac12$.

\begin{theorem}[Near-critical amplitude]
\label{a:thm:nearcritical}
As $p\uparrow\tfrac12$, with $\varepsilon = 1-2p\downarrow0$,
\begin{equation}
\label{a:eq:Anearcritical}
A(p) = 2\varepsilon\lrbs{1+O\lrb{\varepsilon\log\frac1\varepsilon}} .
\end{equation}
In particular
\begin{equation}
\label{a:eq:Aasym}
A(p)\sim2(1-2p),
\qquad
\frac{1}{A(p)}\sim\frac{1}{2(1-2p)} .
\end{equation}
\end{theorem}

\begin{proof}
Write $r = 1-\varepsilon$. The identity
\begin{equation}
\label{a:eq:denomsplit}
\frac{1}{2-S_n} = \frac12 + \frac{S_n}{2(2-S_n)}
\end{equation}
turns the exact series \cref{a:eq:Aseries} into
\begin{equation}
\label{a:eq:seriesdecomp}
\frac{1}{A(p)} = 1 + \frac{1}{2\varepsilon} + D_\varepsilon,
\qquad
D_\varepsilon := \sum_{n=0}^{\infty}r^n\,\frac{S_n}{2(2-S_n)} ,
\end{equation}
the geometric sum $\sum_nr^n = 1/\varepsilon$ having been evaluated. Only a
logarithmic bound on the remainder $D_\varepsilon$ is needed.

The map $s\mapsto ps(2-s)$ is increasing in $p$ and in $s$ on $s\in[0,1]$, so
comparing \cref{a:eq:SnIter} at parameter $p$ with the same recursion at $p=\tfrac12$
from the common initial value $S_0=1$ gives, by induction,
$S_n(p)\le S_n(\tfrac12)$ for every $n$. At criticality the reciprocal increments
satisfy
\begin{equation}
\label{a:eq:recipincrement}
\frac{1}{S_{n+1}(\tfrac12)}-\frac{1}{S_n(\tfrac12)}
= \frac{1}{2\bigl(1-S_n(\tfrac12)/2\bigr)}\;\ge\;\frac12 ,
\end{equation}
so $1/S_n(\tfrac12)\ge1+n/2$ and hence $S_n(\tfrac12)\le2/(n+2)$. Since
$2-S_n\ge1$, these two facts combine to give
\begin{align}
0\le D_\varepsilon
&\le \frac12\sum_{n=0}^{\infty}r^nS_n(p)
 \;\le\; \frac12\sum_{n=0}^{\infty}r^n\,\frac{2}{n+2}
 \;=\; \sum_{n=0}^{\infty}\frac{r^n}{n+2}
\notag\\
&= \frac{-\log(1-r)-r}{r^2}
 \;=\; O\lrb{\log\frac1\varepsilon} .
\label{a:eq:Dbound}
\end{align}
Factoring $1/(2\varepsilon)$ out of \cref{a:eq:seriesdecomp},
\begin{equation*}
\frac{1}{A(p)} = \frac{1}{2\varepsilon}
\Bigl[1 + 2\varepsilon\bigl(1+D_\varepsilon\bigr)\Bigr],
\end{equation*}
and the quantity following the leading $1$ is $O(\varepsilon\log(1/\varepsilon))$.
Inverting gives \cref{a:eq:Anearcritical}.
\end{proof}

The proof locates the logarithm precisely: it comes from the near-critical decay
$S_n\lesssim2/(n+2)$ persisting across the whole range $n\lesssim1/\varepsilon$ over
which the geometric factor $r^n$ still carries weight. The
summand $r^n/(2-S_n)$ involves two decays on very different scales: the geometric
factor $r^n = (1-\varepsilon)^n$ decays on the scale $n\sim1/\varepsilon$, which is
large, whereas $S_n$ is close to its critical behaviour $2/n$ and has already
fallen to near zero after $O(1)$ generations. For the overwhelming majority of the
terms that carry weight, therefore, $S_n\approx0$ and the summand is approximately
$r^n/2$; summing gives $1/A\approx1/(2\varepsilon)$ at once. The rigorous argument
above is that estimate with the neglected mass bounded rather than discarded.

\begin{remark}[The critical gradient]
\label{a:rem:gradient}
By \cref{a:eq:Aasym}, $A(p)$ meets the axis at $p=\tfrac12$ with gradient $-4$:
\begin{equation}
\label{a:eq:gradient}
A(p) = 2-4p+o(1-2p)\qquad(p\to\tfrac12^-).
\end{equation}
This is the single most useful diagnostic in \cref{a:sec:GA}: since $1/A$ diverges
as $A\to0$, an error in $A$ of vanishing absolute size still produces an unbounded
error in the quasi-stationary mean, so a candidate accurate to plotting resolution
across the whole interval can fail here, and fail fatally.
\end{remark}

\begin{remark}[Rate of approach]
\label{a:rem:rate}
Numerically, $A(p)/2\varepsilon$ approaches $1$ rather slowly. Evaluating
\cref{a:eq:prodFormula} in high precision gives $A/2\varepsilon = 0.90987$,
$0.98612$, $0.99814$, $0.99977$ and $0.999972$ at $\varepsilon=2\times10^{-2}$ down
to $2\times10^{-6}$. An exploratory fit to the deficit gives
\begin{equation}
\label{a:eq:empiricalcorrection}
1-\frac{A}{2\varepsilon}\approx\varepsilon\lrb{\log\frac1\varepsilon+c},
\qquad c\approx0.78,
\end{equation}
consistent with \cref{a:eq:Anearcritical}. The constant $c$ is empirical.
\end{remark}

\subsection*{Why the continuum route is worse}

The same leading asymptotic can be reached through the Riccati approximation of
\ChM. The exact difference is
\begin{equation*}
S_{n+1}-S_n = (2p-1)S_n - pS_n^2
= -\varepsilon S_n - \tfrac12S_n^2 + \tfrac{\varepsilon}{2}S_n^2 ,
\end{equation*}
and dropping the $O(\varepsilon S_n^2)$ term gives
$\dd S/\dd n = -\varepsilon S-\tfrac12S^2$. Setting $u=1/S$ linearises this to
$u' = \varepsilon u+\tfrac12$, with solution
\begin{equation*}
u(n) = C\ee^{\varepsilon n} - \frac{1}{2\varepsilon},
\qquad
S(n) = \frac{1}{C\ee^{\varepsilon n}-\frac{1}{2\varepsilon}} .
\end{equation*}
Matching $S_n\sim A\ee^{-\varepsilon n}$ for small $\varepsilon$ identifies
$A=1/C$, and matching to the critical decay $S_n\sim2/n$ in the appropriate double
limit fixes $C\sim1/(2\varepsilon)$, recovering \cref{a:eq:Aasym}. Both steps are
heuristic and the second is delicate; worse, the route gives no control at all on
the error, and so cannot see the logarithm of \cref{a:eq:Anearcritical}.
